\section{Discussion}
\label{sec:discussion}

\paragraph{What a practitioner should take from this.} Three things, in decreasing order of how
confident we are in them. Running longer does not cost reproducibility: the spread across seeds is
flat in the number of updates, so there is no reason to cut a run short to keep results stable.
Batch composition is the dominant lever, and it acts as $P^{-1/2}$, so the returns are real but slow
-- quadrupling the batch halves the spread. And the step size matters less than one would expect,
entering as $D^{+1/2}$, which means a run can be made faster without being made much less
reproducible, up to the saturation boundary below.

\paragraph{On reporting.} The honest consequence of Section~\ref{sec:exp-seeds} is not that error
bars are now free. It is that they are cheap in a way they were not: the scaling transfers, so one
calibration -- a handful of seeds at one setting -- lets a practitioner carry an error bar across
changes of batch and step size, which is where most of the variation between reported configurations
comes from. What we could not deliver is the absolute level from a single run, and
Section~\ref{sec:exp-seeds} says exactly why the obvious route fails.

\paragraph{Against drift targeting, in one case.} Section~\ref{sec:drift} recommends specifying a
run by a drift target rather than a learning rate, and Section~\ref{sec:experiments} shows that
target transferring across batch size where a learning rate does not. Section~\ref{sec:exp-seeds}
shows the same technique destroying two runs. The mechanism is specific and worth stating plainly:
$\E_x[p(1-p)] \to 0$ as a task is solved, the signal goes with it, and a controller holding drift
constant answers by raising the step without bound. Any drift-targeted schedule needs a floor on the
signal or a cap on the step; ours had neither.

\paragraph{Limits.} The linearisation behind Proposition~\ref{prop:stationary} assumes the two runs
stay close enough that the mean gradient field is locally linear between them, which is the same
assumption that fails at saturation. The largest model measured is $0.5$B parameters, and the tasks
are short-answer and exactly verifiable; whether the same picture holds for reasoning-length
responses, where a single rollout carries far more entropy, is the obvious next measurement and we
have not made it. The contraction time is treated as a single number and is not one. And the
estimator family of \eqref{eq:family} excludes clipping, which makes the update a nonlinear function
of the batch and breaks the linearity the noise decomposition relies on.

\paragraph{What we would do next.} Measure the spectrum of the contraction rather than a single
timescale, since that is what stands between the scaling law and an absolute error bar. Repeat the
measurement where responses are long, which is both the regime that matters most and the one where
the within-prompt noise term should be largest. And check whether the same balance explains a
phenomenon we did not set out to study: if divergence between two seeds is stationary, then so is
divergence from any perturbation, including a change of data order or of hardware numerics, and the
same instrument should predict those.
